
\begin{derivation}[从\eqnrefs{eq:schwinger-parameter-r9}到\eqnrefs{eq:feynman-two-general-r68}]
两个分母的 Feynman 参数并不是独立记忆公式，它直接来自两个 Schwinger 参数的重新参数化。分别对 $A^{\nu_A}$ 和 $B^{\nu_B}$ 使用\eqnrefs{eq:schwinger-parameter-r9}，得到
\begin{align*}
 \frac1{A^{\nu_A}B^{\nu_B}}
 &=\frac1{\Gamma(\nu_A)\Gamma(\nu_B)}
 \int_0^\infty\dd\tau_1\dd\tau_2\,
 \tau_1^{\nu_A-1}\tau_2^{\nu_B-1}
 \ee^{-A\tau_1-B\tau_2}.
\end{align*}
令
\begin{equation*}
 T=\tau_1+\tau_2,
 \qquad x=\frac{\tau_1}{T},
 \qquad
 \tau_1=xT,
 \quad \tau_2=(1-x)T.
\end{equation*}
Jacobian 为 $T$，于是
\begin{align*}
 \frac1{A^{\nu_A}B^{\nu_B}}
 &=\frac1{\Gamma(\nu_A)\Gamma(\nu_B)}
 \int_0^1\dd x\,x^{\nu_A-1}(1-x)^{\nu_B-1}\\
 &\quad\times\int_0^\infty\dd T\,
 T^{\nu_A+\nu_B-1}
 \ee^{-T[xA+(1-x)B]}.
\end{align*}
最后一个积分再次是 Gamma 函数，直接给出正文\eqnrefs{eq:feynman-two-general-r68}。取 $\nu_A=\nu_B=1$ 即恢复\eqnrefs{eq:feynman-parameter-proof-r9}。
\end{derivation}

\begin{derivation}[从\eqnrefs{eq:schwinger-parameter-r9}到\eqnrefs{eq:feynman-three-denominator-r14}]
三个分母只比上一个推导多一个比例变量。对 $A,B,C$ 各使用一次幂的 Schwinger 表示，
\begin{equation*}
 \frac1{ABC}
 =\int_0^\infty\dd\tau_1\dd\tau_2\dd\tau_3\,
 \ee^{-A\tau_1-B\tau_2-C\tau_3}.
\end{equation*}
定义
\begin{equation*}
 T=\tau_1+\tau_2+\tau_3,
 \qquad x=\frac{\tau_1}{T},
 \quad y=\frac{\tau_2}{T},
 \quad z=1-x-y,
\end{equation*}
则 $x,y,z\ge0$、$x+y+z=1$，而 Jacobian 为 $T^2$。因此
\begin{align*}
 \frac1{ABC}
 &=\int_0^1\dd x\int_0^{1-x}\dd y
 \int_0^\infty\dd T\,T^2
 \ee^{-T[xA+yB+(1-x-y)C]}\\
 &=2\int_0^1\dd x\int_0^{1-x}\dd y\,
 \frac1{[xA+yB+(1-x-y)C]^3},
\end{align*}
其中用了 $\int_0^\infty T^2\ee^{-DT}\dd T=\Gamma(3)/D^3=2/D^3$。把三角区域写成带 $\delta(x+y+z-1)$ 的对称单纯形测度，就得到正文\eqnrefs{eq:feynman-three-denominator-r14}。
\end{derivation}







\begin{derivation}[从\eqnrefs{eq:euclidean-master-schwinger-dp68}到\eqnrefs{eq:euclidean-master-integral-dp21}]
在正文给出的绝对收敛域内，可以交换 Schwinger 参数与圈动量积分。将\eqnrefs{eq:euclidean-master-schwinger-dp68} 代入，
\begin{align*}
 I_{d,\nu}(\Delta)
 &=\frac1{\Gamma(\nu)}
 \int_0^\infty\dd s\,s^{\nu-1}\ee^{-s\Delta}
 \int\frac{\dd^d \ell_E}{(2\pi)^d}\ee^{-s\ell_E^2}.
\end{align*}
用\eqnrefs{eq:d-dimensional-gaussian-r9} 完成圈动量积分，
\begin{align*}
 I_{d,\nu}(\Delta)
 &=\frac1{(4\pi)^{d/2}\Gamma(\nu)}
 \int_0^\infty\dd s\,
 s^{\nu-d/2-1}\ee^{-s\Delta}.
\end{align*}
再令 $u=s\Delta$，得到
\begin{equation*}
 \int_0^\infty\dd s\,s^{\nu-d/2-1}\ee^{-s\Delta}
 =\Delta^{d/2-\nu}\Gamma(\nu-d/2),
\end{equation*}
从而得到正文\eqnrefs{eq:euclidean-master-integral-dp21}。
\end{derivation}



\begin{derivation}[从\eqnrefs{eq:dimreg-universal-factor-dp68}到\eqnrefs{eq:dimreg-universal-expansion-dp68}]
维数正规化中的方案常数来自三个普通小参数展开。首先由 $\Gamma(1+\epsilon)=\epsilon\Gamma(\epsilon)$ 和
\begin{equation*}
 \Gamma(1+\epsilon)=1-\gamma_E\epsilon+\order(\epsilon^2)
\end{equation*}
得到
\begin{equation*}
 \Gamma(\epsilon)=\frac1\epsilon-\gamma_E+\order(\epsilon).
\end{equation*}
另外
\begin{align*}
 (4\pi)^\epsilon
 &=1+\epsilon\ln4\pi+\order(\epsilon^2),\\
 \left(\frac{\Delta}{\mu^2}\right)^{-\epsilon}
 &=1-\epsilon\ln\frac{\Delta}{\mu^2}+\order(\epsilon^2).
\end{align*}
三者相乘并保留有限阶，得到
\begin{equation*}
 \mathcal U_\epsilon(\Delta)
 =\frac1\epsilon-\gamma_E+\ln4\pi
 -\ln\frac{\Delta}{\mu^2}+\order(\epsilon),
\end{equation*}
即正文\eqnrefs{eq:dimreg-universal-expansion-dp68}。$\overline{\rm MS}$ 方案定义为在这个通用展开中一并减去 $1/\epsilon-\gamma_E+\ln4\pi$；这一步定义了 $\overline{\rm MS}$ 的减除方案；母积分本身仍由前面的维数正则化恒等式给出。
\end{derivation}
